\chapter{Fichiers produits et reproduction}
\label{ann:reproduction}

\section{Organisation des entrées et des notebooks}
Les 24 CSV sources sont organisés dans \code{data/raw/2019} à \code{data/raw/2024}.
Chaque dossier contient \code{caracteristiques.csv}, \code{lieux.csv},
\code{vehicules.csv} et \code{usagers.csv}. Le README du projet décrit la récupération
des sources. Les données volumineuses sont exclues du dépôt Git et des sources du rapport.

\begin{table}[htbp]\centering\small\renewcommand{\arraystretch}{1.18}
\caption{Dépendances entre les étapes du projet.}
\begin{tabularx}{\textwidth}{@{}lLL@{}}\toprule
Notebook & Entrées & Sorties principales\\\midrule
01 & Quatre CSV de 2024 & Tableaux de découverte dans le notebook.\\
02 & Les 24 CSV & Audit et règles de préparation.\\
03 & Les 24 CSV & Tables préparées, schéma, manifeste et six bilans.\\
04 & Table usagers et schéma du 03 & Dix tableaux, huit figures et métadonnées.\\
05 & Table usagers et schéma du 03 & Sept tableaux de bilan, réplications, quatre figures et métadonnées.\\
06 & Table usagers, schéma, bilan du 03 et module ML & Comparaison sur 2023, configuration, diagnostics et trois figures.\\
07 & Entrées du 06 et configuration retenue & Vérification de reconstruction, test 2024, diagnostics et trois figures.\\\bottomrule
\end{tabularx}
\end{table}

Chaque notebook recharge ses entrées depuis un espace de variables neuf. Les notebooks
06 et 07 utilisent le module \code{src/baac_ml.py} ; l'encodeur reste appris dans la pipeline.

\section{Installer et lancer l'étude}
Avec Python 3.12 et GNU Make disponibles sous Linux ou macOS, depuis la racine du projet~:
\begin{verbatim}
make install
.venv/bin/python -m jupyter lab
\end{verbatim}
Le fichier \code{requirements.txt} fixe les versions des bibliothèques directement
utilisées. Leurs dépendances sont installées par \code{pip} selon les contraintes
 déclarées. Le noyau Jupyter choisi doit utiliser ce venv. Les sept notebooks sont
 ensuite exécutés dans l'ordre, chacun depuis un noyau neuf.

La commande \code{make clean} supprime uniquement \code{.venv}, après \code{deactivate}
si nécessaire. \code{make install} le recrée. Le Makefile n'exécute pas les analyses.

\section{Lire les sorties et retrouver leur provenance}
Dans \code{data/processed}, les tables principales sont
\path{baac_usagers_metropole_2019_2024.csv.gz} et
\path{baac_lieux_sources_metropole_2019_2024.csv.gz}. La compression gzip ne change
pas leur unité d'observation. Le schéma \path{baac_regles_et_schema.json}, version
 de règles 1.1, fournit les types et le marqueur \code{__BAAC_MANQUANT__}.
Le manifeste \path{baac_sources_manifest.json} contient les empreintes des sources.
Les six bilans documentent effectifs, géographie, identifiants, jointures, lieux et nettoyage.

Les figures et tableaux sont classés par étape dans les sous-dossiers
\code{04_exploration}, \code{05_statistics}, \code{06_modeling} et \code{07_evaluation},
sous \code{reports/figures} et \code{reports/tables}. Une ligne du fichier
\path{05_replicats_contrastes.csv.gz} représente un résultat de rééchantillonnage,
et non un usager~: ses 16\,000 lignes correspondent à 2\,000 réplications, quatre
contrastes et deux événements.

Le notebook 07 lit \path{06_configuration_retenue.json}.
La copie \path{07_configuration_evaluee.json} assure la traçabilité du test.
Les JSON conservent les règles et empreintes ; les CSV donnent les populations,
scores et diagnostics.

\section{Contrôles de reproduction réalisés}
La revue finale du 17 octobre 2026 a installé les dépendances du projet dans un venv
neuf. La commande \code{pip check} n'a détecté aucun conflit. Les 105 cellules Python
 des sept notebooks ont ensuite été exécutées, sans modification, dans sept processus
neufs, en repartant des 24 CSV sans table préparée ni résultat préexistant.

Les contrôles internes passent, les deux tables compressées reconstruites sont
identiques octet par octet à celles fournies et les résultats statistiques sont
reproduits. Les 18 figures sont régénérées. La logistique retient à nouveau $C=1$~;
les écarts des scores de validation et de test restent inférieurs à $10^{-4}$.

Les chiffres publiés proviennent des exports sous Python 3.12.3. La revue utilise
Python 3.12.14 et les mêmes versions directes des bibliothèques. Son environnement est
décrit dans \path{resultats/environnement_revue.json}. Les noyaux Jupyter sont bloqués
par les restrictions de sockets : les cellules sont exécutées directement en Python,
sans prétendre valider l'interface Jupyter sur la machine de l'auteur.

\section{Modifier et compiler le rapport}
Les sources du rapport contiennent les 18 figures scientifiques, le schéma relationnel
 dessiné en LaTeX, les résultats agrégés nécessaires et un script qui génère 17 tableaux
LaTeX. Le script utilise uniquement la bibliothèque standard Python. Il ne réexécute
ni le nettoyage, ni le bootstrap, ni les modèles.

Depuis le dossier contenant \code{main.tex}~:
\begin{verbatim}
python3 scripts/generer_tableaux.py
pdflatex main.tex
\end{verbatim}
Répéter pdfLaTeX trois fois pour résoudre les renvois et le sommaire. Le README
précise les prérequis et l'importation dans Overleaf. Les chapitres sont séparés
pour faciliter les corrections.
\FloatBarrier
